\section{Curvature on real two-planes}\label{sec:curvature}

We retain the base potential and Hartogs metric of
Proposition~\ref{metric:transfer}.  Section~\ref{sec:pinching} will prove
the two uniform sectional-curvature bounds.  Here we compute the full
curvature tensor and reduce its value on an arbitrary real two-plane
to a scalar expression.  In particular, the calculation includes real
planes that are not complex lines.

Throughout this section, constants in tensor estimates depend only on
the uniform centered-chart bounds for $h=(\phi_{i\bar j})$ in
Proposition~\ref{metric:transfer}; the base dimension is two.
We take $\lambda\ge1$.  The base form remains
$I=(\psi_{i\bar j})$, with $\psi=-\log(1-|z|^2)$.

\subsection{A pointwise holomorphic gauge and the tensor blocks}

Fix a point with $w\ne0$, choose centered ball coordinates at its base
point, and choose a local branch of $\log w$.  Put
\[
 s=\log t=\phi+\log w+\overline{\log w},\qquad
 x=\frac{t}{1-t},\qquad q=x(1+x).
\]
In these coordinates the centered base point is zero.  Remove the first
and pure second derivatives of the weight by the holomorphic fiber change
\[
 P(z)=\sum_i\phi_i(0)z_i+
       \frac12\sum_{i,k}\phi_{ik}(0)z_i z_k,
 \qquad v=\log w+P(z).
\]
More precisely, $s=v+\bar v+\widetilde\phi$, where
$\widetilde\phi=\phi-P-\bar P$.  At the point under consideration,
\begin{equation}\label{curv:gauge}
 s_i=s_{ik}=0,\qquad s_v=s_{\bar v}=1,\qquad
 s_{i\bar j}=h_{i\bar j}.
\end{equation}
All higher derivatives of $s$ involving a vertical index vanish.
Subtracting $P+\bar P$ changes neither $h$ nor its component
derivatives in the centered base coordinates.

The metric at the point consequently has blocks
\[
 \widehat g=\begin{pmatrix}G&0\\0&q\end{pmatrix},
 \qquad G=\lambda I+xh.
\]
At the center, $I$ is the identity matrix.  Normalize the vertical
vector to $q^{-1/2}\partial_v$ and denote its index by $0$.
Horizontal indices $i,j,k,l$ range over $1,2$.
In this pointwise frame the metric is $G+V$, where $V$ is the
Hermitian form whose only nonzero entry is $V_{0\bar0}=1$.
Horizontal forms are extended by zero on the vertical direction.
For any Hermitian form $P$, define
\begin{equation}\label{curv:symmetric-tensor}
 (S_P)_{a\bar b c\bar d}
   =P_{a\bar b}P_{c\bar d}+P_{a\bar d}P_{c\bar b}.
\end{equation}

Our complex curvature components use the convention
\begin{equation}\label{curv:coordinate-formula}
 R_{a\bar b c\bar d}
 =-\mathcal K_{ac\bar b\bar d}
   +\widehat g^{p\bar e}
       \mathcal K_{ac\bar e}\mathcal K_{p\bar b\bar d}.
\end{equation}
Repeated indices in this formula are summed over all coordinates;
repeated horizontal indices below are summed over $1,2$.
The inverse coefficients satisfy
$\sum_p\widehat g^{p\bar e}\widehat g_{p\bar j}
=\delta_{ej}$, so their numerical array is the transpose of the
ordinary inverse of the displayed metric matrix.
We relate \eqref{curv:coordinate-formula} to real sectional curvature
below, with the fixed convention $g=2\operatorname{Re}\widehat g$.

\begin{proposition}\label{curv:tensor}
In the frame just described, define the horizontal Hermitian form
\[
 D=(1+2x)h-qhG^{-1}h.
\]
The full curvature tensor is
\begin{equation}\label{curv:tensor-identity}
 R=-\lambda S_I-qS_h-S_V
    -(S_{D+V}-S_D-S_V)+T_4+T_3,
\end{equation}
where $T_j$ has exactly $j$ horizontal indices.  Its nonzero
components, up to the K\"ahler symmetries and conjugation, are
\begin{align}
 (T_4)_{i\bar j k\bar l}
   &=-x h_{i\bar j,k\bar l}
      +x^2G^{p\bar e}h_{i\bar e,k}h_{p\bar j,\bar l},
       \label{curv:four-horizontal-error}\\
 (T_3)_{i\bar j k\bar0}
   &=\sqrt q\bigl(-h_{i\bar j,k}
           +xG^{p\bar e}h_{i\bar e,k}h_{p\bar j}\bigr).
       \label{curv:three-horizontal-error}
\end{align}
Measured by the Euclidean component norms in this frame, uniformly
for $\lambda\ge1$,
\begin{equation}\label{curv:component-errors}
 \|T_4\|\le C_4x,\qquad \|T_3\|\le C_3\sqrt q.
\end{equation}
If $\lambda$ is at least the uniform upper bound for the eigenvalues
of $h$ at centered-chart origins, then
\begin{equation}\label{curv:D-bounds}
 (1+x)h\le D\le(1+2x)h.
\end{equation}
\end{proposition}

\begin{proof}
Write $F(s)=-\log(1-e^s)$.  Direct differentiation gives
\[
 F'=x,\qquad F''=q,\qquad F'''=q(1+2x),\qquad
 F''''=q(1+6x+6x^2).
\]
The centered-chart expression for $\psi$ differs from
$-\log(1-|z|^2)$ by a pluriharmonic function and a constant.
Its derivatives of type $(2,1)$ therefore vanish at the center,
and its derivatives of type $(2,2)$ are $S_I$.

For brevity in the derivative tables, put
\[
 A_{ik\bar j}=h_{i\bar j,k},\qquad
 B_{ik\bar j\bar l}=h_{i\bar j,k\bar l}.
\]
Before normalizing the vertical vector, the only possibly nonzero
third derivatives of type $(2,1)$, up to symmetry, are
\begin{equation}\label{curv:third-derivatives}
 \begin{aligned}
 \mathcal K_{ik\bar j}&=xA_{ik\bar j},\qquad
 \mathcal K_{iv\bar j}=qh_{i\bar j},\qquad
 \mathcal K_{vv\bar v}=q(1+2x).
 \end{aligned}
\end{equation}
The fourth derivatives of type $(2,2)$, up to symmetry and
conjugation, are
\begin{equation}\label{curv:fourth-derivatives}
 \begin{aligned}
 \mathcal K_{ik\bar j\bar l}
   &=\lambda(S_I)_{i\bar j k\bar l}
       +xB_{ik\bar j\bar l}+q(S_h)_{i\bar j k\bar l},\\
 \mathcal K_{ik\bar j\bar v}&=qA_{ik\bar j},\\
 \mathcal K_{iv\bar j\bar v}&=q(1+2x)h_{i\bar j},\\
 \mathcal K_{vv\bar v\bar v}&=q(1+6x+6x^2).
 \end{aligned}
\end{equation}
To check that these lists are complete, apply the chain rule to
$F(s)$ and use \eqref{curv:gauge}.  In the four-horizontal
derivative, the two pairings of mixed second derivatives of $s$
give $qS_h$; the third pairing contains
$s_{ik}s_{\bar j\bar l}=0$.  All other terms contain a vanished
first or pure second horizontal derivative.  In particular,
$\mathcal K_{ik\bar v}=\mathcal K_{vv\bar j}=0$,
$\mathcal K_{ik\bar v\bar v}=0$, and the fourth derivatives
with one horizontal and three vertical indices vanish.
No vanishing of pure third horizontal derivatives is needed.

Substitution into \eqref{curv:coordinate-formula} gives
\eqref{curv:four-horizontal-error} and
\eqref{curv:three-horizontal-error}.  For clarity, all remaining
curvature types in the normalized frame are
\begin{equation}\label{curv:remaining-blocks}
 \begin{aligned}
 R_{i\bar j0\bar0}
   &=-(1+2x)h_{i\bar j}
       +q(hG^{-1}h)_{i\bar j}=-D_{i\bar j},\\
 R_{i\bar0k\bar0}&=0,\qquad
 R_{i\bar00\bar0}=0,\qquad
 R_{0\bar00\bar0}=-2.
 \end{aligned}
\end{equation}
The last identity follows because its unnormalized value is
\[
 -q(1+6x+6x^2)+q(1+2x)^2=-2q^2,
\]
and the four vertical normalization factors multiply to $q^{-2}$.
The tensors $S_V$ and $S_{D+V}-S_D-S_V$ give precisely the
vertical and balanced two-horizontal blocks in
\eqref{curv:remaining-blocks}.  The K\"ahler symmetries and
conjugation now account for every component, proving
\eqref{curv:tensor-identity}.

The lower bound $h\ge cI$ implies
\[
 \|G^{-1}\|\le\frac1{\lambda+cx},\qquad
 x^2\|G^{-1}\|\le\frac{x}{c}.
\]
The assumed component-derivative bounds for $h$, applied to
\eqref{curv:four-horizontal-error} and
\eqref{curv:three-horizontal-error}, prove
\eqref{curv:component-errors} with constants independent of
$\lambda\ge1$ and $x>0$.

Finally, $G=\lambda I+xh$ commutes with $h$ at the center.
On an eigendirection of $h$ with eigenvalue $d$, the ratio of $D$
to $h$ is
\[
 1+2x-\frac{x(1+x)d}{\lambda+xd}.
\]
The subtracted fraction is nonnegative, and it is at most $x$
when $\lambda\ge d$.  This proves \eqref{curv:D-bounds}.
\end{proof}

Although the logarithmic fiber coordinate was used only for $w\ne0$,
the tensor at the zero section can be checked directly in regular
coordinates.  A holomorphic fiber change $W=e^{P(z)+a}w$, with
constant $a$ and a polynomial of degree at most two, makes the new
weight have value zero and vanishing first and pure second
derivatives at the base point.  Its mixed Hessian remains $h$.
Through total degree four, the terms relevant to the curvature are
\[
 \lambda\psi+
 |W|^2\left(1+\sum_{i,j}h_{i\bar j}z_i\bar z_j\right)
 +\frac12|W|^4.
\]
At the origin the metric is $\operatorname{diag}(\lambda I,1)$,
all relevant third derivatives vanish, and direct differentiation
gives
\begin{equation}\label{curv:zero-section}
 R_{i\bar j k\bar l}=-\lambda(S_I)_{i\bar j k\bar l},
 \qquad R_{i\bar j0\bar0}=-h_{i\bar j},
 \qquad R_{0\bar00\bar0}=-2,
\end{equation}
with all other types zero.  In the original fiber coordinate the
normalized vertical vector is $e^{-\phi/2}\partial_w$.
Thus the formulas above hold at the zero section as well by setting
$x=q=0$, $G=\lambda I$, $D=h$, and $T_3=T_4=0$, and using this
regular normalized vertical vector.

\subsection{Real curvature and Gram determinants}

Let $U_1,U_2$ be real linearly independent tangent vectors and let
$u_1,u_2$ be their $(1,0)$ parts, so $U_j=u_j+\bar u_j$.
For a Hermitian form $P$, linear in its first argument, define
\begin{equation}\label{curv:gram-quantities}
 \begin{aligned}
 a_P&=P(u_1,u_1)P(u_2,u_2)
             -(\operatorname{Re}P(u_1,u_2))^2,\\
 b_P&=\operatorname{Im}P(u_1,u_2),\qquad
 m_P=a_{P+V}-a_P-a_V\quad\text{when $P$ is horizontal}.
 \end{aligned}
\end{equation}
Thus $a_P$ is the determinant of the real Gram matrix for
$\operatorname{Re}P$, $b_P$ is the imaginary Hermitian pairing, and
$m_P$ is the cross term in that determinant after adding $V$.
For the metric form $\widehat g$, the real squared area is
$4a_{\widehat g}$ because $g=2\operatorname{Re}\widehat g$.
For a tensor $T$ with the K\"ahler symmetries, also put
\[
 \mathcal H_T=T_{u_1\bar u_1u_2\bar u_2}
       -\operatorname{Re}T_{u_1\bar u_2u_1\bar u_2}.
\]
Here contraction with complex vectors has its usual multilinear
meaning.  These definitions allow the horizontal or vertical projection
of the real plane to be degenerate.

\begin{lemma}\label{curv:real-plane}
For the metric convention $g=2\operatorname{Re}\widehat g$ and
the curvature convention \eqref{curv:coordinate-formula},
\begin{equation}\label{curv:sectional}
 K_g(\operatorname{span}_{\mathbb R}\{U_1,U_2\})
       =\frac{\mathcal H_R}{2a_{\widehat g}}.
\end{equation}
Define the matrix
\begin{equation}\label{curv:xi-definition}
 \Xi^{a\bar b}=u_1^a\overline{u_2^b}
                   -u_2^a\overline{u_1^b}.
\end{equation}
For every tensor $T$ with the K\"ahler symmetries and Hermitian
conjugation symmetry,
\begin{equation}\label{curv:xi-contraction}
 \mathcal H_T=-\frac12
 T_{a\bar b c\bar d}\Xi^{a\bar b}\Xi^{c\bar d}.
\end{equation}
For every Hermitian form $P$,
\begin{equation}\label{curv:S-contraction}
 \mathcal H_{S_P}=a_P+3b_P^2.
\end{equation}
\end{lemma}

\begin{proof}
The real squared area of the pair is $4a_{\widehat g}$, because
$g(U_i,U_j)=2\operatorname{Re}\widehat g(u_i,u_j)$.
In particular $a_{\widehat g}>0$ for a real linearly independent
pair.  Write
\[
 A=R_{u_1\bar u_1u_2\bar u_2},\qquad
 B=R_{u_1\bar u_2u_1\bar u_2}.
\]
The K\"ahler symmetries in the expansion on
$U_j=u_j+\bar u_j$ give the real numerator
\[
 R^{\mathbb R}(U_1,U_2,U_2,U_1)
      =A-B-\bar B+A=2\mathcal H_R.
\]
Division by the real squared area proves \eqref{curv:sectional}.
Expanding the right-hand side of \eqref{curv:xi-contraction}
gives the same expression $A-\operatorname{Re}B$ for $T$.
Finally, if $c_P=P(u_1,u_2)$, direct contraction of
\eqref{curv:symmetric-tensor} gives
\[
 \mathcal H_{S_P}
   =P(u_1,u_1)P(u_2,u_2)+|c_P|^2-2\operatorname{Re}(c_P^2)
   =a_P+3b_P^2.
\]
\end{proof}

For example, the disk potential $-\log(1-|w|^2)$ has normalized
complex curvature component $-2$ by
\eqref{curv:remaining-blocks}.  Taking $u_2=iu_1$ in
\eqref{curv:sectional} gives real sectional curvature $-2$.
This fixes the normalization also on complex lines.

\begin{lemma}\label{curv:gram}
For a positive semidefinite Hermitian form $P$,
$a_P\ge b_P^2$.  For the vertical form $V$ one has $a_V=b_V^2$.
If $P$ is horizontal and positive semidefinite, $m_P$ is nonnegative,
linear and monotone in $P$, and
\begin{equation}\label{curv:mixed-gram}
 m_P\ge2\sqrt{a_Pa_V}.
\end{equation}
When $a_Pa_V>0$, equality in \eqref{curv:mixed-gram} holds exactly
when the two real Gram matrices of the pair for $P$ and $V$ are
proportional.
\end{lemma}

\begin{proof}
Complex Cauchy--Schwarz gives
$P(u_1,u_1)P(u_2,u_2)\ge|P(u_1,u_2)|^2$, proving
$a_P\ge b_P^2$.  Equality holds for $V$ because its complex rank
is one.
Let $A$ and $B$ be the two-by-two real Gram matrices with entries
$\operatorname{Re}P(u_i,u_j)$ and
$\operatorname{Re}V(u_i,u_j)$, respectively.  Then
\[
 a_P=\det A,\qquad a_V=\det B,\qquad
 m_P=\operatorname{tr}(\operatorname{adj}(A)B).
\]
For positive definite $A$, put $C=A^{-1/2}BA^{-1/2}$.
The last expression equals $(\det A)\operatorname{tr}C$,
whereas $\sqrt{a_Pa_V}=(\det A)\sqrt{\det C}$.
The inequality $\operatorname{tr}C\ge2\sqrt{\det C}$ proves
\eqref{curv:mixed-gram}, and continuity gives the semidefinite
case.  If both determinants are positive, equality means that
$C$ is a scalar matrix, which is exactly the claimed proportionality.
The displayed formula for $m_P$ is linear in the entries of $P$.
Its nonnegativity for positive semidefinite $P$ then also proves
monotonicity.
\end{proof}

\subsection{Error terms and the curvature numerator}

Decompose $u_j=(u_{j,H},u_{j,0})$ in the pointwise frame.
Write $\Xi_{HH}$ for the horizontal two-by-two block of $\Xi$
and $\Xi_{H\bar0}$ for its horizontal-to-vertical column.
The other mixed block is the negative conjugate transpose of this
column.  Write $\Xi_{\mathrm{mix}}$ for the matrix consisting of
both mixed blocks, with all other entries zero.
All norms of these blocks below are Euclidean Frobenius norms.

Direct expansion of \eqref{curv:xi-definition} gives
\begin{equation}\label{curv:xi-norms}
 \begin{aligned}
 \|\Xi_{HH}\|^2&=2(a_I+b_I^2)\le4a_I,\\
 \|\Xi_{H\bar0}\|^2&=m_I+2b_Ib_V,\qquad
 \|\Xi_{\mathrm{mix}}\|^2=2(m_I+2b_Ib_V)\le4m_I.
 \end{aligned}
\end{equation}
For the second identity, explicitly expand
$\|u_{1,H}\overline{u_{2,0}}
      -u_{2,H}\overline{u_{1,0}}\|^2$.
The last inequality follows from
$|b_Ib_V|\le\sqrt{a_Ia_V}\le m_I/2$,
by Lemma~\ref{curv:gram}.

\begin{lemma}\label{curv:errors}
For every real linearly independent tangent pair,
\begin{equation}\label{curv:contracted-errors}
 |\mathcal H_{T_4}|\le Cx a_I,\qquad
 |\mathcal H_{T_3}|\le C\sqrt q\,\sqrt{a_I m_I}.
\end{equation}
More explicitly, the three-horizontal error satisfies
\begin{equation}\label{curv:mixed-error}
 |\mathcal H_{T_3}|
   \le C\sqrt q\,\|\Xi_{HH}\|\,
                         \|\Xi_{\mathrm{mix}}\|.
\end{equation}
\end{lemma}

\begin{proof}
In \eqref{curv:xi-contraction}, the tensor $T_4$ pairs two
horizontal blocks of $\Xi$.  The tensor $T_3$ pairs one horizontal
block and one mixed block, since it has exactly three horizontal
indices.  Cauchy--Schwarz and \eqref{curv:component-errors} give
\eqref{curv:mixed-error} and
$|\mathcal H_{T_4}|\le Cx\|\Xi_{HH}\|^2$.
Now use \eqref{curv:xi-norms}.  These estimates remain valid when
one of the projected real areas vanishes.
\end{proof}

We can now express the entire real curvature numerator.
By \eqref{curv:S-contraction} and $a_V=b_V^2$,
\[
 \mathcal H_{S_V}=4a_V,\qquad
 \mathcal H_{S_{D+V}-S_D-S_V}=m_D+6b_Db_V.
\]
Consequently \eqref{curv:tensor-identity} gives
\begin{equation}\label{curv:negative-numerator}
 \begin{aligned}
 -\mathcal H_R={}&\lambda(a_I+3b_I^2)+q(a_h+3b_h^2)
        +4a_V+m_D+6b_Db_V\\
       &-\mathcal H_{T_4}-\mathcal H_{T_3}.
 \end{aligned}
\end{equation}
Its denominator in \eqref{curv:sectional} is positive for every real
two-plane.  It remains to prove that
\eqref{curv:negative-numerator} is uniformly positive relative to
that denominator for sufficiently large fixed $\lambda$, and that
the quotient is also uniformly bounded above.  The mixed term
$6b_Db_V$ can have either sign; Section~\ref{sec:pinching} treats
it together with the possible degeneration of the projected planes.
